% =====================================================================
%  BOZZA - Sommario (italiano)
%
%  Versione italiana qualitativa dell'abstract, allineata alla nuova
%  struttura a tre capitoli scritti (Introduzione, Contributo,
%  Implementazione). Niente numeri specifici: i risultati quantitativi
%  arriveranno nel capitolo dedicato.
% =====================================================================
\chapter*{Sommario}
\addcontentsline{toc}{chapter}{Sommario}

Nuotatori, snorkeler, sub e kayaker sono poco visibili alle imbarcazioni che
condividono le acque costiere. Le tecnologie di sicurezza pensate per le grandi
navi, come l'AIS, i beacon satellitari di emergenza e il tracciamento cellulare,
sono troppo costose, troppo lente o legate a un'infrastruttura che al largo non
esiste. Un'alternativa praticabile \`e usare la radio Wi-Fi gi\`a presente in
molti dispositivi come beacon di prossimit\`a senza infrastruttura: delle boe
intelligenti trasmettono periodicamente brevi frame broadcast con la propria
posizione, che le imbarcazioni vicine ascoltano.

Questa tesi lavora su un simulatore a eventi discreti per il beaconing broadcast
IEEE~802.11 e si concentra sulla propagazione della consapevolezza oltre il
singolo hop radio. Il contributo principale \`e l'implementazione e il confronto
di due meccanismi che estendono la conoscenza che ogni nodo ha della rete: la
modalit\`a \emph{append}, che aggiunge a ogni beacon la lista dei nodi gi\`a
conosciuti senza inviare frame in pi\`u, e la modalit\`a \emph{forwarded}, che
rilancia attivamente i beacon ricevuti entro un limite di hop. Accanto a questi
meccanismi si descrivono i protocolli che decidono \emph{quando} una boa
trasmette: il beaconing statico a intervallo fisso (SBP), l'adattamento basato
sulla densit\`a (ADAB) e lo scheduler context-aware ACAB, che combina densit\`a
dei vicini, tempo dall'ultimo contatto e mobilit\`a del nodo.

Il modello di canale \`e calibrato su misure in mare aperto, la mobilit\`a segue
il modello Random Waypoint e la popolazione di boe varia nel tempo per imitare gli
utenti che entrano ed escono dall'acqua. I capitoli che seguono presentano la
motivazione del lavoro, il progetto dei due meccanismi multihop con le formule che
ne governano il comportamento, e l'implementazione nel simulatore. La valutazione
sperimentale e le conclusioni sono lasciate ai capitoli successivi.
